\section{A survey of the geometry behind the conjecture}\label{sec:survey}

This section explains the objects and the obstruction measured by the conjecture, with particular attention to the different scopes of the two detection arguments. Readers familiar with Schubert varieties and local divisor theory may proceed to Section~\ref{sec:introduction}. The proofs in the rest of the paper do not depend on the survey.

The motivating question concerns singular spaces with finitely many natural strata: does the ambient local geometry at a general point of each maximal singular component determine where a given algebraic defect occurs, including on smaller strata? Here the defect is failure to be Gorenstein. On a Schubert variety this means that the canonical sheaf is not a line bundle. Divisor class groups encode that failure, and localization supplies the maps by which it is tested.

\subsection{Flags, Schubert cells, and Bruhat order}
A complete flag in $\C^n$ is a chain $0=F_0\subset F_1\subset\cdots\subset F_n=\C^n$ with $\dim F_i=i$. These flags form the smooth projective variety $\GL_n/B$, where $B$ is the group of invertible upper triangular matrices. Fix the standard flag $E_j=\langle e_1,\ldots,e_j\rangle$. The $B$-orbits record relative positions with respect to $E_\bullet$ and are the Schubert cells \cite[Section~1.2]{Bri05}.

The cell $C_w$ indexed by $w\in S_n$ contains the coordinate flag $F_i=\langle e_{w(1)},\ldots,e_{w(i)}\rangle$. Row elimination identifies it with $\A^{\ell(w)}$, where $\ell(w)$ counts inversions. Its closure $X_w=\overline{C_w}$ is a Schubert variety. Explicitly, it is defined by
\[
 \dim(F_i\cap E_j)\ge\#\{a\le i:w(a)\le j\}.
\]
In matrix charts these incidence inequalities become rank bounds, and hence equations given by minors \cite[Section~3]{WY08}. The cells are smooth, while their closures may be singular.

Bruhat order is defined by $u\le w$ if $C_u\subseteq X_w$, equivalently $X_u\subseteq X_w$; thus $X_w=\coprod_{u\le w}C_u$. A cover $u\lessdot w$ has $\ell(w)=\ell(u)+1$. In $S_2$, the identity labels a point and the other permutation labels the open affine cell in $\PP^1$. In general the longest element labels the whole flag variety. Bruhat order is a partial order: a dimension inequality alone does not imply a closure inclusion \cite[Chapter~2]{BB05}.

A partial flag remembers only specified subspaces; a Grassmannian remembers one. Forgetting steps gives $G/B\to G/P$, a locally trivial projection with smooth flag variety fibre. This allows local singularity questions to pass between complete and partial flags; Section~4 gives the compatible detector maps. Symplectic and orthogonal groups give flags compatible with bilinear forms. Their Weyl groups replace ordinary permutations. Types~B and~C have the same abstract signed-permutation Weyl group but different root lengths, which produce different integral character matrices below.

\subsection{What a local ring and a generic point remember}
For $X=\Spec R$, a point is a prime ideal $\mathfrak p$ and its local ring is $R_{\mathfrak p}$. At a closed complex point this ring consists of rational functions defined on some neighbourhood. General prime ideals are equally important \cite[Chapter~II, Sections~2--3]{Har77}.

An irreducible closed subvariety $Z$ has a generic point $\eta_Z$ with closure $Z$. The ambient ring $\OO_{X,\eta_Z}$ records transverse behaviour along a nonempty open subset of $Z$: its dimension is $\codim_X Z$, while its residue field is $\C(Z)$.

\Needspace{5\baselineskip}
For example, for the line $Z=V(y)$ in $\A^2$, one has
\[
 \OO_{\A^2,\eta_Z}=\C[x,y]_{(y)}\simeq\C(x)[y]_{(y)},
\]
a discrete valuation ring with parameter $y$ and residue field $\C(x)$. In contrast, the intrinsic ring $\OO_{Z,\eta_Z}$ is the field $\C(x)$. Every generic-point test here concerns the ambient ring.

If $x\in Z$, then $\eta_Z$ is a generalization of $x$ and there is a localization map $\OO_{X,x}\to\OO_{X,\eta_Z}$. Freeness of a module persists under this localization. Its converse generally fails: a property may hold at general points of a component and fail at special ones. The conjecture asks for a specific converse for a particular obstruction on Schubert varieties.

This also explains the field extensions used in Section~4. A generic point along a positive-dimensional cell has a function field as residue field, even though the variety is defined over $\C$. The local arguments therefore retain an arbitrary extension $F\supset\C$ and use the same $F$ on every closed subslice participating in one argument.

\subsection{Regularity, normality, and Gorensteinness}
A Noetherian local ring $(R,\mathfrak m)$ is regular if $\dim_{R/\mathfrak m}(\mathfrak m/\mathfrak m^2)=\dim R$. For varieties over $\C$, regularity is the local algebraic form of smoothness. A singular point may nevertheless satisfy weaker conditions that control duality and divisors.

An integral variety is normal if its local rings are integrally closed domains. In codimension one these rings are discrete valuation rings, so rational functions have integral orders of vanishing along prime divisors. Normality also permits functions and reflexive sheaves to be recovered across complements of codimension at least two \cite[Chapter~II, Section~6]{Har77}\cite[Tags 031S, 0EBJ]{The26}.

A local ring is Cohen--Macaulay if its depth equals its dimension. Depth is the maximal length of a regular sequence in the maximal ideal, each element being a nonzerodivisor after the preceding ones have been set equal to zero. Complete intersections are Cohen--Macaulay, but the latter class is much larger \cite[Chapters~1--2]{BH93}.

The canonical module expresses local duality. On a smooth $d$-dimensional variety, the canonical sheaf is the line bundle $\bigwedge^d\Omega_X^1$. On a normal variety it is its rank-one reflexive extension from the smooth locus. For the Cohen--Macaulay local rings considered here, Gorensteinness is equivalent to the canonical module being free of rank one \cite[Chapter~3]{BH93}\cite[Tag 0DW9]{The26}. Thus a singular ring can retain duality with one local generator. The standard implications \cite[Tag 0AWV]{The26} are
\[
 \text{regular}\ \Longrightarrow\ \text{local complete intersection}
 \ \Longrightarrow\ \text{Gorenstein}
 \ \Longrightarrow\ \text{Cohen--Macaulay}.
\]
Normality is a separate condition. Schubert varieties in this paper are normal and Cohen--Macaulay \cite{Ram85,Bri05}. Our question is where their canonical sheaves stop being line bundles, not merely where the varieties become singular.

\subsection{Divisors and the group containing the obstruction}
A prime Weil divisor on a normal integral variety is an irreducible closed subvariety of codimension one. Weil divisors are finite integer combinations of these; a rational function $f\ne0$ defines $\operatorname{div}(f)=\sum_E\operatorname{ord}_E(f)E$. Divisors modulo principal divisors form $\Cl(X)$. A Cartier divisor is locally the divisor of one nonzero rational function, with unit ratios on overlaps, and determines a line bundle. Every Weil divisor on a smooth variety is Cartier, whereas this can fail on a normal singular variety \cite[Chapter~II, Section~6]{Har77}\cite{Fos73}.

The local group $\Cl(\OO_{X,x})$ measures the failure near $x$: $D$ is Cartier at $x$ precisely when its local class is zero. The qualification ``local'' matters. A point on $\PP^1$ defines the globally nontrivial line bundle $\OO_{\PP^1}(1)$, although its class is zero in every local class group. Globally, $\Cl(\PP^1)\simeq\Z$.

Equivalently, a Weil class is a rank-one reflexive sheaf $\mathcal F$, with $\mathcal F\simeq\mathcal F^{**}$ and $\mathcal F^*=\mathcal Hom(\mathcal F,\OO_X)$. The group operation is tensor product followed by double dual \cite[Tags 0EBJ, 0EBM]{The26}. Cartierness at a point is freeness of rank one over the local ring. This description supplies a closed nonfree locus; the flat descent used below is proved in Lemma~\ref{alg:flat}.

Let $K_X$ represent the canonical sheaf. Since $X$ is Cohen--Macaulay, it is Gorenstein at $x$ if and only if $[K_X]=0$ in $\Cl(\OO_{X,x})$.
A normal local ring is factorial precisely when its class group is zero. Factoriality therefore tests every class; Gorensteinness tests the distinguished canonical class. A divisor is $\Q$-Cartier at $x$ if some positive multiple is Cartier, meaning that its local class is torsion. Its Cartier index is the order of that class. The analogous condition for $K_X$ is $\Q$-Gorensteinness. Nonzero finite-order and infinite-order canonical classes both obstruct Gorensteinness, but a torsion-detection theorem applies only to the former.

\subsection{Three model calculations}
\begin{example}[A singular Gorenstein point with infinite class group]\label{ex:node}
The rank-at-most-one $2\times2$ matrix variety $D_{2,2}=\Spec\C[a,b,c,d]/(ad-bc)$ has dimension three. The gradient $(d,-c,-b,a)$ vanishes only at its vertex, where the tangent space has dimension four. It is the ordinary threefold node; being a hypersurface, it is Gorenstein.

The ruling primes $P=(a,b)$ and $Q=(a,c)$ satisfy $\operatorname{div}(a)=P+Q$. Inverting $a$ gives the factorial ring $\C[a^{\pm1},b,c]$, whose units are $\C^*a^{\Z}$. Divisor localization therefore identifies the global class group with
\[
 (\Z[P]\oplus\Z[Q])/\Z(P+Q)\simeq\Z.
\]
Its generator stays nonzero in the vertex local ring: if a homogeneous symbolic power became principal there, graded Nakayama would make it principal globally. Section~3 proves this comparison in detail; see also \cite[Section~8]{BV88}. Consequently the local class of $a_0P+b_0Q$ is $(a_0-b_0)[P]$. A singular Gorenstein point can carry many non-Cartier Weil divisors, and comparing the two ruling coefficients is an integral test.
\end{example}

\begin{example}[Finite-order divisors and the canonical class]\label{ex:veronese}
The invariant ring for simultaneous negation on $\A^2$ is
\[
 R=\C[s^2,st,t^2]\simeq\C[u,v,w]/(uw-v^2).
\]
It is normal and Gorenstein. At its vertex the prime $P=(u,v)$ is not principal, while $\operatorname{div}(u)=2P$ and $R[u^{-1}]=\C[u^{\pm1},v]$. Thus the vertex local class group is $\Z/2$: there is a torsion divisor obstruction, although the canonical class is zero.

For simultaneous negation on $\A^3$, the invariant ring is $\C[x_ix_j:1\le i\le j\le3]$. Its vertex local class group is $\Z/2$ \cite[Section~1.2, Proposition~I]{Wat74}. The involution acts on $dx_1\wedge dx_2\wedge dx_3$ by $-1$; the group has no pseudo-reflections and is not contained in $\SL_3$, so Watanabe's local criterion gives non-Gorensteinness \cite[Theorem~2]{Wat74}. The canonical class corresponds to the determinant character \cite[Section~2, Theorem~$1^\prime$, pp.~383--384]{Wat74}. It is the nonzero class, and its square is trivial, giving canonical index two. Equivalently, the toric lattice is the even-sum sublattice of $\Z^3$; the canonical divisor has odd total parity. These scalar Veronese cones are the type~C models studied in Section~5.
\end{example}

\begin{example}[An infinite-order canonical obstruction]\label{ex:determinantal}
For $r,s\ge2$, the rank-at-most-one variety $D_{r,s}$ of $r\times s$ matrices is normal and Cohen--Macaulay, has dimension $r+s-1$, and has vertex local class group $\Z$. Its row and column rulings represent primitive opposite classes. Up to the choice of sign for the generator, its canonical class is $(r-s)$ times that generator \cite[Corollary~8.4, Theorem~8.8 and Corollary~8.9]{BV88}. Indeed, the variety is the cone over $\PP^{r-1}\times\PP^{s-1}$ embedded by $\OO(1,1)$; the canonical bundle $\OO(-r,-s)$ maps to the difference of its coefficients after quotienting by the diagonal hyperplane class. The homogeneous local comparison is the same as in Example~\ref{ex:node}. Thus $D_{2,3}$ has a nonzero canonical class of infinite order and is not $\Q$-Gorenstein at its vertex.
\end{example}

These models have isolated singularities. Their punctured spectra are smooth, so Weil classes are line bundles on the puncture. For a nonisolated singularity a reflexive module need not be invertible on the whole puncture. The torsion argument first passes to the generic point of a component of the nonfree locus; only then does it obtain the necessary line bundle.

\subsection{From global classification to a locus problem}
A permutation contains a pattern $\pi\in S_m$ if some $m$ entries have the same relative order as $\pi$. Lakshmibai--Sandhya proved that a type~A Schubert variety is smooth precisely when its permutation avoids $3412$ and $4231$ \cite{LS90}\cite[Theorem~7.6]{WY23}. This converts a geometric property into a finite combinatorial test, but it does not determine Gorensteinness: Example~\ref{ex:node} is singular and Gorenstein.

The singular locus is closed and $B$-stable, hence a finite union of Schubert subvarieties. Its irreducible components are called maximal singular components, with maximality understood by inclusion, not dimension. Their generic singularities in type~A are described in \cite{BW03,Cor03,KLR03,Man01b,Man01a}. They supply a manageable list of ambient local models. The remaining issue is whether those models detect all failures on deeper strata.

Woo--Yong's global classification determines which complete type~A Schubert varieties are Gorenstein everywhere \cite{WY06}. Their Corollary~1 says that Gorensteinness at every maximal singular generic point implies Gorensteinness everywhere. The locus conjecture asks more: every non-Gorenstein point should lie on a maximal singular component that is already generically non-Gorenstein. With $\M(X)=\Irr(\Sing X)$, the asserted identity is
\begin{equation}\label{eq:survey-wy}
 \NGor(X)=\bigcup_{\substack{Z\in\M(X)\\\OO_{X,\eta_Z}\text{ not Gorenstein}}}Z.
\end{equation}
Openness gives the right-to-left inclusion. The other inclusion is substantive, particularly when generically Gorenstein and non-Gorenstein components coexist. Removing the known bad components does not permit another application of a theorem about complete Schubert varieties.

Subsequent work includes the interval-pattern approach \cite{WY08}, Perrin's determination of the Gorenstein locus in minuscule Schubert varieties \cite[Theorem~0.3 and Corollary~0.4]{Per09}, boundary blow-ups and a Demazure-subword criterion \cite{DS19}, and the survey \cite{WY23}. Global Cartier-class results \cite{LRY25} and canonical trace calculations for Grassmannian Schubert cycles and determinantal rings \cite{Kim26} address related questions. Section~\ref{sec:introduction} specifies their scopes.

\subsection{Why the conjecture leads to divisor detection}
At $x\in X$, localizing at the maximal singular generic points whose closures contain $x$ gives
\begin{equation}\label{eq:survey-rho}
 \rho_x:\Cl(\OO_{X,x})\longrightarrow
 \bigoplus_{\substack{Z\in\M(X)\\x\in Z}}\Cl(\OO_{X,\eta_Z}).
\end{equation}
The containment is necessary for these localization maps to exist. The canonical form of the conjecture asks whether $\rho_x([K_X])=0$ implies $[K_X]=0$. Injectivity of the whole map is stronger. Theorem~\ref{thm:mainA} proves that stronger assertion in type~A, including infinite-order classes, and consequently also determines the nonfactorial locus \cite[Conjecture~6.14]{WY08}\cite[Conjecture~7.39]{WY23}.

Theorem~\ref{thm:maintorsion} for arbitrary finite type asserts injectivity only on the source torsion subgroup. This qualification cannot be dropped: even the abstract map $\Z\to0$ is injective on its zero torsion subgroup while killing every nonzero integer. Likewise, vanishing of generic images does not make a local canonical class finite-order.

For a class already known to have finite order, injection preserves its order. The order in a finite direct sum is the least common multiple of the component orders, so a divisor known to be $\Q$-Cartier at $x$ satisfies the index formula of Corollary~\ref{cor:indices}. An empty least common multiple is one. The source hypothesis supplies finite order; the generic tests determine that finite order exactly.

\subsection{Transverse slices remove smooth directions}
Let $B^-$ be the opposite Borel and $C^u=B^-uB/B$. The Kazhdan--Lusztig slice $N_{u,w}=X_w\cap C^u$ has distinguished point $u$ and dimension $\ell(w)-\ell(u)$. There is an actual open product
\[
 X_w\cap\Omega_u\simeq C_u\times N_{u,w}
 \simeq\A^{\ell(u)}\times N_{u,w},
\]
compatible with Schubert strata \cite[Lemma~3.2]{WY08}\cite[Lemma~2.1]{KWY13}. The affine factor records movement along the cell, and the slice records transverse local geometry. Freeness of the relevant reflexive modules can be compared through the faithfully flat product.

In type~A, free matrix entries have root characters. Bruhat rank inequalities determine which root directions survive in the tangent space. A contracting positive grading then shows that the surviving entries generate the actual coordinate ring. This is why matrix algebra and permutation inequalities enter the proof together.

A tangent space records only first-order information, and a tangent cone is an associated graded object. Neither determines the original singularity: for instance, $y^2=x^3$ and $y^2=x^5$ have the same tangent cone $y^2=0$. Here the product and the determinantal models are actual scheme isomorphisms, carrying the prime divisors and their multiplicities. An unmarked generic-singularity classification alone would not identify which two divisors become the matrix rulings. Theorem~\ref{ta:filtered} supplies that placement information.

\subsection{The type A proof: comparing coefficients}
The node calculation illustrates the mechanism: when a matrix entry has two reduced ruling components, a divisor supported on them is locally trivial precisely when its two coefficients agree.

Choose the first-difference coordinate $t$ at $v<w$. Its nonvanishing locus is $N_{v,w}[t^{-1}]\simeq\Gm\times N_{u,w}$ with $\ell(u)=\ell(v)+1$. Induction handles divisor classes on this lower-dimensional slice.

A class vanishing on the open set can then be written $[D]=\sum_i c_i[E_i]$, where the $E_i$ are the prime components of $t=0$ and the $c_i$ are integers. This representation need not be unique. Reducedness gives the principal relation $\sum_i[E_i]=0$, so equality of all the coefficients suffices.

Some components come from Schubert boundary divisors and are called vertical. A possible additional component meets the dense cell and is called horizontal. For neighbouring active vertical components, a maximal element of a filtered Bruhat intersection gives a rank-one matrix cone in which the two chosen divisors are the rulings of one entry. The cone is singular only at its vertex, so its label corresponds to a maximal singular component of the ambient variety, an allowed detector.

The horizontal component requires an additional construction. Proposition~\ref{ta:horizontal} moves to a larger stratum while retaining a non-Cartier vertical divisor. A global divisor-class relation identifies the horizontal contribution in another two-ruling comparison. Thus every required neighbouring difference, including the last horizontal comparison, is detected. If every detector vanishes, all $c_i$ coincide and the principal relation gives $[D]=0$.

The proof therefore needs integral coefficients, reduced multiplicities, actual divisor markings, and globally maximal singular components. These properties are established separately; they are not inferred from a tangent-cone resemblance.

\subsection{The general-type proof: finite covers and integral degrees}
Finite-order classes admit an additional interpretation. If a line bundle $L$ on $U$ has a chosen trivialization $L^{\otimes m}\simeq\OO_U$, the powers of $L$ form a cyclic finite algebra. Locally its cover has equation $z^m=a$ with $a$ a unit, so in characteristic zero it is finite étale. Kummer theory relates these covers to torsion line bundles, with unit ambiguity $\Gamma(U,\OO_U)^*/\Gamma(U,\OO_U)^{*m}$ \cite{Gro71}\cite[Tag 03PK]{The26}.

To apply this, one first takes the generic point of an irreducible component of the nonfree locus. The reflexive module is then free on the punctured local spectrum. Passing to a strict henselization, a henselian local ring with separably closed residue field, allows simple roots to lift \cite[Tag 04GE]{The26}. Finite étale covers of the entire strict local spectrum split, although covers of its puncture may not. In the normal strictly henselian local rings of dimension at least two used here, units on the puncture extend and have $m$th roots. The unit ambiguity disappears. The descent argument uses injectivity for the relevant punctured Picard group; it does not identify the entire ordinary and strict local class groups.

The first detection method restricts covers to the Schubert boundary. Its top pieces correspond to coatoms $v\lessdot w$, and smaller Bruhat labels describe their overlaps. For $d=\ell(w)-\ell(u)$, the order complex of $(u,w)$ is a sphere $S^{d-2}$ \cite{BW82}. When $d\ge4$, its simple connectivity makes labels for split sheets on the pieces compatible. A local Lefschetz theorem then lifts the splitting from the boundary to the puncture \cite[Tag 0BLJ]{The26}. The restriction functor is used only for full faithfulness.

The second method uses a regular proper birational source, obtained by restricting a Bott--Samelson resolution to the slice. This restricted source is not itself asserted to be a tower of projective lines. Let $E_j$ be all of its prime divisors lying over the vertex. For a Kummer equation $z^m=a$, valuations off that fibre are divisible by $m$, so
\[
 \operatorname{div}(a)=mD+\sum_jc_jE_j.
\]
Taking degrees on actual proper integral cycles $Z_i$ in the fibre gives
\[
 N(c_j)_j\equiv0\pmod m,
 \qquad N_{ij}=\deg\bigl(\OO(E_j)|_{Z_i}\bigr).
\]
Smith normal form gives an annihilator for the torsion class. In particular, a full-column primitive matrix eliminates it: an integer left inverse exists, or equivalently the full-column maximal minors have gcd one. Rational rank is insufficient; the matrix $(2)$ permits a nonzero solution modulo $2$.

The geometry identifies every divisor column and realizes the character rows as integer combinations of genuine curves, using the Bott--Samelson fibre geometry of \cite{Ha04}. The effective Levi reduction bounds the required rank by slice dimension. In type~C the high-dimensional induction identifies a slice with nonzero strict puncture torsion as an actual scalar Veronese cone $\A^d/\{\pm1\}$. Its singularity is isolated, so it cannot hide torsion at a nonmaximal singular label. The bounded type~$F_4$ calculation and the type~$G_2$ word argument complete the finite-type result. The finite computations establish only their stated bounded inputs; geometry supplies the all-ranks and all-fields conclusions.

\subsection{The type D family and infinite-order obstruction}
For $r\ge k\ge3$, Section~\ref{sec:type-d} constructs a Schubert variety $X\subset\OG(k,2r+4)$ with a unique maximal singular component $S$ of codimension three. The transverse model at $\eta_S$ is the Gorenstein node. Nevertheless the ambient variety is non-$\Q$-Gorenstein at every point of a smaller locus $T\simeq\Gr(k,r)\subset S$, of codimension $2k$ in $X$. The Grassmannian $T$ is itself smooth; it is the ambient local rings of $X$ along $T$ that fail the canonical condition.

An incidence variety $I$ is a normal Gorenstein determinantal hypersurface and maps properly and birationally to $X$. It is isomorphic to $X$ away from $T$ and has fibre $\PP^1$ over every point of $T$, with no exceptional prime divisor. Along a contracted fibre curve $C$, the canonical line bundle of $I$ has degree $2k-4>0$.

If $mK_X$ were Cartier near its image, its pullback would have degree zero on $C$. Since there are no exceptional prime divisors, comparing canonical Weil divisors identifies that pullback with $mK_I$, whose degree is $m(2k-4)$. This contradiction applies to every scheme point of $T$, and proves an infinite-order local canonical obstruction.

The smallest case $(r,k)=(3,3)$ has type~$D_5$, dimension six, and $T$ a point. The case $(5,4)$ has type~$D_7$, dimension twelve, and $T\simeq\PP^4$. These orthogonal Grassmannians are nonminuscule, so the examples do not conflict with Perrin's theorem. They explain the additional term $\NonQGor(X)$ in Corollary~\ref{cor:indices}.

\subsection{Reading routes and background}
The three main assertions have distinct scopes: all divisor classes in type~A, finite-order classes in arbitrary finite type, and an infinite-order canonical obstruction in type~D. For the original conjecture, Sections~3 and~6 give the coefficient argument and its consequences. For the local-cover and lattice argument, Sections~4 and~5 give the geometry and induction. Section~7 is independent of the finite certificates and can be read directly after the divisor background.

For further background, \cite{Bri05} treats flag varieties and their divisor theory, while \cite{WY23} develops the combinatorial perspective. Canonical modules are treated in \cite{BH93}, divisor class groups in \cite{Har77,Fos73}, and determinantal models in \cite{BV88}. Bruhat topology is covered in \cite{BB05,BW82}, finite covers in \cite{Gro71}, and the normal affine semigroup methods in \cite{CLS11}. The new arguments begin with the marked detectors and their assembly into the local statements.
